% Generated draft: requires mathematical and source review.
\documentclass[11pt]{article}
\usepackage[margin=25mm]{geometry}
\usepackage{fontspec}
\setmainfont{Latin Modern Roman}
\usepackage{amsmath,amssymb,mathtools,unicode-math}
\setmathfont{Latin Modern Math}
\usepackage{enumitem,needspace,xurl,hyperref}
\setlength{\parindent}{0pt}
\setlength{\parskip}{5pt}
\setlength{\emergencystretch}{3em}
\newcommand{\N}{\mathbb{N}}
\newcommand{\Z}{\mathbb{Z}}
\newcommand{\Q}{\mathbb{Q}}
\newcommand{\R}{\mathbb{R}}
\newcommand{\C}{\mathbb{C}}
\newcommand{\F}{\mathbb{F}}
\newcommand{\A}{\mathbb{A}}
\newcommand{\PP}{\mathbb{P}}
\newcommand{\E}{\mathbb{E}}
\newcommand{\eps}{\varepsilon}
\newcommand{\dd}{\,\mathrm{d}}
\begin{document}
\section*{UnsolvedMath ID: 800455}
\textbf{Source problem code:} catalogue entry
\subsection{Definitions and mathematical statement}
For a positive integer \(n\), let \([n]=\{1,\ldots,n\}\). For a real-valued function \(f\) on \(\{-1,1\}^n\), use its unique multilinear representation
\[
f(x)=\sum_{S\subseteq[n]}\widehat f(S)\prod_{i\in S}x_i
\]
and define \(\deg(f)=\max\{|S|:\widehat f(S)\ne0\}\) for nonzero \(f\). This is the standard Boolean-cube interpretation of degree; the supplied source does not explicitly define that term.

Let \(\mathcal A_n\) consist of all functions
\[
f:\{-1,1\}^n\longrightarrow[-1,-2/3]\cup[2/3,1]
\]
satisfying, for every \(x\in\{-1,1\}^n\),
\[
\begin{aligned}
\sum_{i=1}^n x_i\ge n/2&\implies f(x)\in[2/3,1],\\
\sum_{i=1}^n x_i\le -n/2&\implies f(x)\in[-1,-2/3].
\end{aligned}
\]
Determine the least possible degree
\[
D(n)=\min_{f\in\mathcal A_n}\deg(f)
\]
as a function of \(n\). In particular, at every input with \(-n/2<\sum_{i=1}^n x_i<n/2\), the value of \(f\) must still lie in \([-1,-2/3]\cup[2/3,1]\), although either sign is permitted. The threshold inequalities include equality. The source does not specify whether an exact formula or an asymptotic characterization of \(D(n)\) is intended.
\subsection{Short English statement}
Determine the minimum multilinear polynomial degree of a real-valued function on the Boolean cube that has magnitude between \(2/3\) and \(1\) everywhere, is positive when the coordinate sum is at least \(n/2\), and is negative when that sum is at most \(-n/2\). Either sign is allowed between the thresholds, but the magnitude restriction remains.
\subsection{Sources}
\begin{itemize}
\item Local UnsolvedMath registry, permanent ID \(800455\), source problem code \texttt{catalogue entry}, title: Analysis of Boolean Functions --- Approximate Degree for Approximate Majority. The supplied registry labels the problem open; that label is not an independently verified mathematical or literature assessment.
\item O'Donnell, \emph{Open Problems in Analysis of Boolean Functions} (2012), Problem 26, ``Approximate Degree for Approximate Majority'': \url{https://arxiv.org/abs/1204.6447}. The supplied record reports access on 2026-10-16 and extraction from arXiv LaTeX; these are supplied provenance statements, not claims of access for this draft.
\item Attribution: the registry lists O'Donnell and 2012 in its proposer fields, while its background gives ``Proposed/Source: Srikanth Srinivasan, 2010.'' Both attributions are retained without resolving their differing roles.
\end{itemize}
\end{document}
